(1) 设 \(y_j=A_j e^{i(kx-\omega t)}\)，得 \(\begin{pmatrix}T_1k^2+\kappa-\rho\omega^2&-\kappa\\-\kappa&T_2k^2+\kappa-\rho\omega^2\end{pmatrix}(A_1,A_2)^T=0\)。非零振幅要求行列式为零。对任意实波数 \(k\)，该矩阵实对称，因而正交的本征模和实频率平方可由能量正定进一步保证。

(2) 解二阶本征值问题得 \(\rho\omega_\pm^2=[(T_1+T_2)k^2/2+\kappa]\pm\sqrt{[(T_1-T_2)k^2/2]^2+\kappa^2}\)。在 \(k=0\)，低支 \(\omega_-(0)=0\) 是两弦同步整体平移，高支 \(\omega_+(0)=\sqrt{2\kappa/\rho}\) 是相对位移的有隙振动。非零耦合 \(\kappa\) 使两支不相交，即使两独立弦在某一波数本可发生频率交叉。

(3) 若 \(T_1=T_2=T\)，本征向量为 \((1,1)\)、\((1,-1)\)，前者连接弹簧始终不形变，对应 \(\omega_-^2=Tk^2/\rho\)，后者把连接弹簧拉伸两倍，对应 \(\omega_+^2=(Tk^2+2\kappa)/\rho\)。当 \(T_1\ne T_2\)，混合角可由 \(\tan2\theta=2\kappa/[(T_1-T_2)k^2]\)（视向量排序可带符号）求得，说明近 \(k=0\) 主要由连接弹簧确定正反相，而大波数主要由各弦张力确定模态归属。

(4) 能量密度 \(\mathcal E=\rho(\dot y_1^2+\dot y_2^2)/2+T_1(y_{1,x})^2/2+T_2(y_{2,x})^2/2+\kappa(y_1-y_2)^2/2\ge0\)。对时间求导并代入动力方程，得到 \(\partial_t\mathcal E+\partial_x\mathcal S=0\)，其中 \(\mathcal S=-T_1\dot y_1y_{1,x}-T_2\dot y_2y_{2,x}\)。无限长或无能流边界下总能量守恒。本征频率的非负性和混合模态的物理一致性因此不依赖对行列式的纯代数判断。
